\chapter*{Úvod}
\addcontentsline{toc}{chapter}{Úvod}

V lednu 2021 podala nizozemská vláda demisi. 
Důvodem byla toeslagenaffaire: automatizovaný systém odhalování podvodů při vyplácení příspěvků na péči o děti nesprávně označil tisíce rodin za podvodníky, přičemž nepřiměřeně zasáhl osoby s dvojím občanstvím. 
Rok předtím Haagský okresní soud zakázal provoz systému SyRI, jenž profiloval příjemce sociálních dávek podle rizika podvodu, s odůvodněním, že porušuje čl. 8 Evropské úmluvy \parencite{RechtbankSyRI20}.
Tyto dva případy jsou pro předkládanou práci výchozí, neboť ukazují dvojí: algoritmické systémy skutečně přijímají rozhodnutí politické povahy, a přitom otázka, kdo za tato rozhodnutí odpovídá, nemá v současné institucionální architektuře jasnou odpověď.


Rozhodnutí, jež bezprostředně určují životní vyhlídky občanů — kdo dostane úvěr, kdo bude pozván k pracovnímu pohovoru, komu bude přiznána sociální dávka, které zprávy se ráno objeví na displeji telefonu, čí hlas bude ve veřejné rozpravě slyšen a čí zůstane ve stínu — jsou dnes ve značné míře přijímána algoritmicky. 
Nejde přitom o rozhodnutí technická v tom smyslu, v jakém je technickým rozhodnutím dimenzování mostní konstrukce. 
Jsou to rozhodnutí o rozdělení zdrojů, příležitostí, viditelnosti a důvěry, tedy rozhodnutí ve vlastním smyslu politická. 
Zvláštností současné situace je, že jsou činěna v prostoru, který nemá ani tvář, ani adresu: nikoli v parlamentu, nýbrž v optimalizační funkci; nikoli po veřejné diskusi, nýbrž před ní, na úrovni technické architektury, k níž veřejnost nemá přístup a jejíž podrobnosti jsou chráněny jako obchodní tajemství.
Kritická literatura posledního desetiletí na tuto situaci odpovídá převážně jazykem škody. Digitální platformy podle ní odcizují pozornost, atrofují schopnost soustředění, uzavírají občany do epistemických bublin, polarizují veřejnou debatu a prostřednictvím mikrotargetingu ovlivňují výsledky voleb. 
Tento jazyk je rétoricky účinný, avšak argumentačně křehký. 
Značná část jeho empirických opor je totiž v odborné literatuře sporná: Rozsáhlá studie Orbenové a Przybylského \parencite{OrbenPrzybylski19} nalezla mezi používáním digitálních technologií a duševní pohodou pouze velmi malou korelaci. Tento výsledek však sám o sobě nevylučuje jiné, heterogenní nebo dlouhodobé účinky digitálního prostředí;
 teze o filtračních bublinách je empiricky opakovaně zpochybňována \parencite{Bruns19,Guess23};
rozsah účinků politického mikrotargetingu zůstává sporný a nelze jej automaticky ztotožnit se změnou volebního chování. Novější experimentální výzkum však ukazuje, že personalizovaná politická sdělení mohou být v určitých podmínkách přesvědčivější než sdělení nepersonalizovaná \parencite{SimchonEdwardsLewandowsky24}.
 Vzniká tak nepříjemná otázka: kdyby se ukázalo, že digitální prostředí polarizuje méně, než se předpokládá, a že jeho vliv na kognitivní schopnosti je zanedbatelný — ztratila by kritika algoritmické moci svůj předmět?
Předkládaná práce tvrdí, že nikoli. 
A právě v tom spočívá její hlavní argumentační záměr.


Ústřední otázka práce zní: jak algoritmicky zprostředkovaná digitální infrastruktura mění podmínky možnosti veřejné racionality — a jaký druh politického problému tato změna představuje?
Otázka se rozpadá do čtyř dílčích:


1. Pojmová: Co je veřejná racionalita a jaké jsou její materiální předpoklady?


2. Diagnostická: Které mechanismy ekonomiky pozornosti a platformového kapitalismu na tyto předpoklady působí a jakým způsobem?


3. Normativní: V čem přesně spočívá politické bezpráví: v porušení autonomie, ve vykořisťování, v epistemické škodě, nebo v něčem jiném?


4. Institucionální: Jaké odpovědi z takto vymezeného bezpráví plynou a proč jsou stávající regulatorní nástroje strukturálně nedostatečné?


Politický problém algoritmické moci nespočívá primárně v měřitelné škodě, kterou digitální technologie působí kognitivním schopnostem jednotlivců, nýbrž ve vzniku nové formy dominance ve smyslu neorepublikánské politické filozofie: nevolené soukromé subjekty disponují svévolnou, neprůhlednou a demokraticky nekontrolovanou mocí nad infrastrukturou, v níž se veřejná racionalita jedině může uskutečňovat.
Kognitivní, epistemické a komunikační škody, jimž je věnována druhá kapitola, nejsou tímto tvrzením popírány. 
Jsou však chápány jako projev problému. Republikánská tradice zde poskytuje rozhodující rozlišení: otrok není svobodný ani při laskavém pánovi, neboť laskavost pána není zaručeným právem, nýbrž libovůlí. 
Analogicky: uživatel algoritmické platformy není svobodný ani tehdy, není-li konkrétní aplikace algoritmu namířena proti němu.
 Dominance nespočívá v uskutečněné škodě, nýbrž v možnosti svévolného zásahu, jež nepodléhá kontrole těch, kdo jsou jí podrobeni.
Tato teze má tři přednosti oproti standardní kritické argumentaci:

Za prvé je robustní vůči empirické revizi: i kdyby se prokázalo, že algoritmická média nezpůsobují měřitelnou kognitivní škodu, politické bezpráví dominance by trvalo.
 

Za druhé přesouvá otázku z terapeutické roviny (jak chránit jednotlivce před závislostí a manipulací) do roviny ústavně-politické (jak rozdělit moc a učinit ji odpovědnou).
  

Za třetí vysvětluje strukturální selhání stávající regulace: GDPR i AI Act vycházejí z individualistického modelu ochrany práv, zatímco dominance je vztah kolektivní a infrastrukturní — a proto individuálním souhlasem neodstranitelný.


Metodou práce je pojmová analýza a kritická rekonstrukce, prováděná v tradici kritické teorie technologie \parencite{Feenberg02} a politické epistemologie \parencite{Coeckelbergh23}.
Práce vychází z výsledků cizích výzkumů a z filozofické literatury.
 Z toho plyne důsledné metodologické opatření, jež zde předesílám: kdekoli se opírám o empirické tvrzení, které je v odborné literatuře sporné, tuto spornost výslovně uvádím a formuluji argument tak, aby na daném tvrzení nezávisel. 
 Sporné empirické teze slouží jako ilustrace mechanismu, nikoli jako premisy hlavního argumentu.


\uv{Umělá inteligence} je v názvu práce pojmem zastřešujícím. 
Skrývají se pod ním tři technicky i politicky odlišné režimy, jejichž ztotožňování je v kritické literatuře běžné a je zdrojem podstatných nepřesností. 
Práce je proto důsledně rozlišuje:


1. Systémy řazení a doporučování (ranking, recommender systems) — organizují viditelnost; působí na veřejnou racionalitu prostřednictvím ekonomiky pozornosti;


2. Prediktivní a skórovací systémy — rozdělují statky a rizika (úvěr, dávka, zaměstnání, míra recidivy); působí prostřednictvím distribuce a rozostření odpovědnosti;


3. Generativní systémy (jazykové a obrazové modely) — produkují obsah; působí prostřednictvím proměny epistemického prostředí.
\uv{Veřejná racionalita} je pojem, který je třeba pečlivě odlišit od Rawlsova \uv{veřejného rozumu} (public reason). 
U Rawlse jde o normativní omezení kladené na obsah ospravedlnění v otázkách základní ústavní úpravy.
 V předkládané práci jde naopak o kolektivní schopnost a o její materiální podmínky: o infrastrukturu, kognitivní předpoklady a sdílený epistemický prostor, bez nichž nemůže být veřejný rozum v Rawlsově smyslu vůbec uplatňován. 
 Veřejná racionalita je tedy podmínkou možnosti veřejného rozumu, nikoli jeho synonymem. 
Systematické vymezení pojmu podává oddíl 2.1.


Otázka po odpovědnosti za jednání, jehož subjekt nelze identifikovat, a po vztahu k entitě, jež napodobuje osobu, aniž by osobou byla, není otázkou novou. 
Židovská a křesťanská tradice ji promýšlela dávno před kybernetikou \parencite{Ellul64}: v podobě golema, jenž jedná, ale nemluví a nenese odpovědnost \parencite{Scholem65}; ve sporu o imago Dei a o hranice lidského tvoření; v Buberově rozlišení mezi vztahem Já–Ty a Já–Ono \parencite{Buber70};
v Levinasově pojetí odpovědnosti zakládané tváří druhého \parencite{Levinas69}; v Jonasově imperativu odpovědnosti za vzdálené důsledky technického jednání \parencite{Jonas84};
 v Arendtově analýze zla, jež nemá zlého původce \parencite{Arendt63,Arendt03}.
Práce tyto zdroje čte jako myšlenkový horizont, z něhož vychází.
 Jejich soustavné rozpracování ve vztahu k algoritmické moci však přesahuje možnosti předkládané práce;
  z tohoto horizontu proto činím produktivním především jediný pojem — Arendtové analýzu zla, jež nemá zlého původce — a tvrdím, že jím lze mezeru odpovědnosti, kterou současná literatura popisuje termíny responsibility gap a antropomorfizace, uchopit přesněji než jazykem soudobé etiky AI.
 Právě v tomto přeznačení spočívá dílčí přínos předkládané práce.


První kapitola rekonstruuje pojmovou genealogii, na jejímž konci stojí ztotožnění myšlení s výpočtem, a ukazuje, jaké politické důsledky z tohoto ztotožnění plynou: možnost delegovat rozhodnutí na systém, jenž nemůže být nositelem odpovědnosti. 
Kapitola dále vymezuje mýtus algoritmické neutrality jako mechanismus depolitizace. 
Zde je třeba předeslat jedno omezení: spor o to, zda výpočetní systémy \uv{rozumějí}, ponechávám otevřený. Ukazuji, že politický argument předkládané práce je na jeho rozhodnutí nezávislý. Pro atribuci odpovědnosti není rozhodující, zda systém rozumí, nýbrž zda existuje identifikovatelný subjekt, jenž za jeho účinky odpovídá.
Druhá kapitola vymezuje pojem veřejné racionality (Kant, Habermas) a analyzuje mechanismy, jimiž ekonomika pozornosti a platformový kapitalismus působí na její materiální předpoklady: kognitivní, epistemické a infrastrukturní. 
Je to kapitola diagnostická, a proto také ta, v níž je nejvíce sporných empirických tvrzení; ta jsou jako sporná označena.
Třetí kapitola provádí normativní a institucionální vyhodnocení. Ukazuje, že technokracie je pokusem přesunout otázky ze sféry komunikativní racionality do sféry instrumentální; ; že algoritmické systémy produkují strukturální mezeru odpovědnosti, kterou lze přesněji popsat pojmem banality zla u Arendtové než jazykem současné etiky AI;
že stávající regulatorní rámce (GDPR, AI Act ve znění tzv. Digital Omnibus z roku 2026) narážejí na individualistický model ochrany práv; a že adekvátní odpovědí je republikánská strategie neutralizace dominance — ochrana, rozdělení moci a budování alternativní infrastruktury.



Práce netvrdí, že digitální technologie jsou samy o sobě škodlivé; problémem je konkrétní institucionální forma jejich organizace, nikoli technologie jako taková. 
Netvrdí, že existoval předdigitální zlatý věk veřejné racionality, k němuž je třeba se vrátit; naopak bere vážně námitku, že nářek nad úpadkem myšlení provázívá každou novou mediální technologii nejméně od Platónova Faidra. 
Nezabývá se debatou o tzv. existenčním riziku pokročilé AI ani vojenským užitím algoritmických systémů. 
Geografickým těžištěm analýzy je euroatlantický prostor; systémy sociálního kreditu v ČLR jsou zmiňovány komparativně, nikoli analyzovány.
